\scanpage{70}
\uuline{Poles de $\mathrm{dn}\,z$} homologues de \uline{$\omega_1$, résidu $i$}, et $-\omega_1$, résidu $-i$

\uuline{\deleted{\unsure{Poles|Zéros}}Zéros de $\mathrm{dn}\,z$} homologues de \uline{$\omega_3$, dérivée $k\,\mathrm{cn}\,\omega_3$}$=i\,\frac{e^{\eta_2\omega_2}}{\sigma^2\omega_2}=-i\,\mathrm{cn}'\omega_2=\pm i\sqrt{1-k^2}$ et de $-\omega_3$, (dérivées symétriques)

Enfin $\mathrm{dn}\,\omega_2=-\mathrm{cn}'\omega_2=\frac{e^{\eta_2\omega_2}}{\sigma^2\omega_2}=\pm\sqrt{1-k^2}\qquad\mathrm{dn}'\omega_2=0$

\uline{formules de dérivation}
\[
\left\{\begin{aligned}
\mathrm{sn}'z&=\mathrm{cn}\,z\,\mathrm{dn}\,z\\
\mathrm{cn}'z&=-\mathrm{sn}\,z\,\mathrm{dn}\,z\\
\mathrm{dn}'z&=-k^2\mathrm{sn}\,z\,\deleted{\unsure{d}}\mathrm{cn}\,z
\end{aligned}\right. \tag*{11}
\]

formules liant $\mathrm{sn}\,z$, $\deleted{\unsure{d}}\mathrm{cn}\,z$, $\mathrm{dn}\,z$ deux à deux
\[
\boxed{\mathrm{cn}^2z+\mathrm{sn}^2z=1\qquad
\deleted{\unsure{k^2}}\mathrm{dn}^2z+k^2\mathrm{sn}^2z=1\qquad
\mathrm{dn}^2z-k^2\mathrm{cn}^2z=1-k^2} \tag*{(12)}
\]

$\mathrm{cn}\,z$ et $\mathrm{dn}\,z$ vérifient des équations analogues à celle de $\mathrm{sn}\,z$
\[
\left\{\begin{aligned}
\mathrm{cn}^{\prime\,2}z&=(1-\mathrm{cn}^2z)(1-k^2\unsure{-|+}k^2\mathrm{cn}^2z)=(1-k^2)(1-\mathrm{cn}^2z)\left(1+\frac{k^2}{1-k^2}\mathrm{cn}^2z\right)\\
\mathrm{dn}^{\prime\,2}z&=(1-\mathrm{dn}^2z)(\mathrm{dn}^2z-(1-k^2)=-(1-k^2)(1-\mathrm{dn}^2z)\left(1-\frac{1}{1-k^2}\mathrm{dn}^2z\right)
\end{aligned}\right. \tag*{(13)}
\]

Aussi $\mathrm{cn}\,z$ et $\mathrm{dn}\,z$ s'expriment\deleted{\unsure{is}}\unsure{-il|-ils} directement au moyen de la fonction $\mathrm{sn}\,z$ qui correspond respectivement à la solution de l'équation $\left(\frac{dZ}{dz}\right)^2=(1-Z^2)(1-{\unsure{\rho|p}}^{2}Z^2)$ \deleted{\unsure{avec|ou|autre}} \unsure{pour|avec} $\rho=\frac{k}{\sqrt{1-k^2}}$ et $\rho=\frac{1}{\sqrt{1-k^2}}$

% archive sheet number: 69
